% ============================================================
% CAPITOLO 27 - SETTIME DI TERZA SPECIE: VII7 MAGGIORE E II7 MINORE
% ============================================================
\chapter{Settime di terza specie: \sigla{VII}{7}{} maggiore e \sigla{II}{7}{} minore}
\label{chap:vii7_ii7}

Lo stesso insieme di note può assumere funzioni diverse secondo il contesto tonale. Due accordi di settima formati dagli stessi suoni (in tonalità relative) hanno ruoli distinti: la \textbf{settima di sensibile} (\sigla{VII}{7}{}) nel modo maggiore e la \textbf{settima sul secondo grado} (\sigla{II}{7}{}) nel modo minore.

\section{La settima di sensibile nel modo maggiore}
\label{sec:vii7_maggiore}

Costruendo un accordo di settima sul VII grado della scala maggiore si ottiene una quadriade.

\begin{definizione}[Settima di sensibile (\sigla{VII}{7}{})]
	\label{def:vii7}
	La \textbf{settima di sensibile} è un accordo di settima costruito sul VII grado della scala maggiore: una triade diminuita (fondamentale, terza minore, quinta diminuita) con l'aggiunta di una settima minore. In Do maggiore: Si-Re-Fa-La.

	È detto anche \textbf{semidiminuito} o \textbf{settima di terza specie} e si indica con \sigla{VIIø}{7}{}.
\end{definizione}

L'accordo contiene tre note della settima di dominante, manca solo la fondamentale. Ha dunque funzione di dominante e risolve sulla tonica.

\begin{regola}[Funzione e risoluzione della \sigla{VII}{7}{}]
	\label{reg:funzione_vii7}
	La \sigla{VII}{7}{} ha funzione di \textbf{dominante}. La risoluzione naturale è sulla tonica (I). Per le regole generali la settima scende di grado (Capitolo \ref{chap:teoria_dissonanza}). Le due note che formano il tritono (Si-Fa in Do maggiore) seguono il loro movimento obbligato:
	\begin{itemize}
		\item la sensibile (Si) sale alla tonica (Do);
		\item la quinta diminuita (Fa) scende alla terza della tonica (Mi);
		\item la settima (La) scende di grado (Sol).
	\end{itemize}
\end{regola}

\begin{esempio}[\sigla{VII}{7}{} in Do maggiore e sua risoluzione]
	\label{ex:vii7_maggiore}
	L'accordo in stato fondamentale risolve sulla tonica. La sensibile Si sale a Do, la settima La e la quinta diminuita Fa scendono di grado (Sol, Mi); il Re della terza scende a Do. La tonica risultante è completa.
	\musica{es_cap13_semidim_do_1}
\end{esempio}

\section{I rivolti della \sigla{VII}{7}{}}
\label{subsec:rivolti_vii7}

Come tutte le settime, anche la \sigla{VII}{7}{} si può rivoltare. Il suo uso più comune è nel \textbf{primo rivolto} (\sigla{VII}{6}{5}).

\begin{regola}[Comportamento delle voci della \sigla{VII}{7}{}]
	\label{reg:comportamento_vii7}
	Alla risoluzione della settima di sensibile valgono queste norme:
	\begin{itemize}
		\item la \textbf{fondamentale} (la sensibile) \textbf{sale di grado} alla tonica e \textbf{non si raddoppia};
		\item la \textbf{quinta} e la \textbf{settima} sono note a risoluzione obbligata: \textbf{scendono di grado} e \textbf{non si raddoppiano};
		\item la \textbf{terza} è parte libera. Quando però si trova nella parte grave, al di sotto della settima, deve \textbf{salire di grado}, per evitare due quinte parallele per moto retto.
	\end{itemize}
\end{regola}

\begin{osservazione}[Disposizione della settima]
	\label{oss:disposizione_settima}
	La settima dell'accordo deve trovarsi a distanza di \textbf{settima} dalla sensibile, mai di seconda o di nona. La distanza di nona è però tollerata al terzo rivolto, purché la settima sia preparata da un accordo consonante.
\end{osservazione}

\begin{center}
\begin{tabularx}{\linewidth}{lLLLL}
\toprule
\textbf{Stato} & \textbf{Cifratura} & \textbf{Grado al basso} & \textbf{Risolve su} & \textbf{Basso obbligato} \\
\midrule
Fond. & 7 & VII & I & sì (sale) \\
I rv. & 6/5 & II & \sigla{I}{6}{} & sì (sale) \\
II rv. & 4/3 & IV & \sigla{I}{6}{} & sì (scende) \\
III rv. & 4/2 & VI & V (\sigla{I}{6}{4}) & sì (scende) \\
\bottomrule
\end{tabularx}
\end{center}

\begin{esempio}[\sigla{VII}{4}{3} e \sigla{VII}{2}{} in Do maggiore]
	\label{ex:vii7_rivolti}
	Il \sigla{VII}{4}{3} (basso Fa) risolve sull'\sigla{I}{6}{} con il basso che scende a Mi; il \sigla{VII}{2}{} (basso La, la settima) risolve sull'\sigla{I}{6}{4} con il basso che scende a Sol. Nel terzo rivolto la settima al basso va preparata (salvo che nella settima diminuita).
	\musica{es_cap13_semidim_do_2}
\end{esempio}

\section{Risoluzioni anticipate della \sigla{VII}{7}{}}
\label{subsec:risoluzioni_anticipate}

Se si hanno a disposizione due movimenti, la settima dell'accordo (in stato fondamentale e nei rivolti) può risolvere \textbf{anticipatamente}: la settima scende di grado e la \sigla{VII}{7}{} diventa direttamente una settima di dominante (\sigla{V}{7}{}) o un suo rivolto, prima di cadere sulla tonica.

\begin{regola}[Risoluzioni anticipate]
	\label{reg:risoluzioni_anticipate}
	\begin{itemize}
		\item $\sigla{VII}{7}{} \rightarrow \sigla{V}{6}{5} \rightarrow \text{I}$
		\item $\sigla{VII}{6}{5} \rightarrow \sigla{V}{4}{3} \rightarrow \text{I}$ (oppure $\sigla{I}{6}{}$)
		\item $\sigla{VII}{4}{3} \rightarrow \sigla{V}{2}{} \rightarrow \sigla{I}{6}{}$
		\item $\sigla{VII}{2}{} \rightarrow \sigla{V}{7}{} \rightarrow \text{I}$
	\end{itemize}
\end{regola}

\begin{esempio}[Risoluzione anticipata \sigla{VII}{7}{} $\rightarrow$ \sigla{V}{6}{5} $\rightarrow$ I in Do maggiore]
	\label{ex:risoluzione_anticipata}
	La settima (La) scende anticipatamente sulla settima di dominante (Sol): il \sigla{VII}{7}{} diventa un \sigla{V}{6}{5} che risolve regolarmente sulla tonica.
	\musica{es_cap12_03}
\end{esempio}

\section{La settima sul secondo grado nel modo minore}
\label{sec:ii7_minore}

Nel modo minore armonico l'accordo di settima costruito sul secondo grado ha la stessa struttura della \sigla{VII}{7}{} del maggiore: una triade diminuita con settima minore.

\begin{definizione}[Settima sul secondo grado (\sigla{II}{7}{})]
	\label{def:ii7_minore}
	La \textbf{settima sul secondo grado} nel modo minore armonico è un accordo semidiminuito con funzione di \textbf{sottodominante}. In La minore: Si-Re-Fa-La, gli stessi suoni della \sigla{VII}{7}{} di Do maggiore. Si indica come \sigla{IIø}{7}{} o semplicemente \sigla{II}{7}{}.
\end{definizione}

Nel relativo maggiore gli stessi suoni formano la settima di sensibile (Definizione \ref{def:vii7}), che porta alla tonica Do. Nel modo minore lo stesso accordo, costruito sul II grado, cambia funzione: è un accordo di sottodominante e, come il \sigla{II}{7}{} del maggiore, risolve sulla dominante (Mi in La minore).

\subsection{La progressione I - \sigla{II}{7}{} - V - I nel modo minore}
\label{subsec:progressione_ii7}

Il \sigla{II}{7}{} del minore si usa soprattutto nella progressione fondamentale tonica - sottodominante - dominante - tonica.

\begin{regola}[Preparazione della settima sul II grado]
	\label{reg:preparazione_ii7}
	La settima del \sigla{II}{7}{}, secondo i principi generali della dissonanza (Capitolo \ref{chap:teoria_dissonanza}), è preparata come consonanza nell'accordo precedente (I). La condotta più fluida si ottiene con il \textbf{primo rivolto} del \sigla{II}{7}{}, con la terza al basso (Re).
\end{regola}

\begin{esempio}[Progressione I - \sigla{II}{6}{5} - V - I in La minore]
	\label{ex:progressione_ii7_minore}
	La settima (La) è preparata nell'accordo di tonica e risolve scendendo di grado (Sol$\sharp$) nell'accordo di dominante.
	\musica{es_cap12_04}
\end{esempio}

\section{Collocazione e rivolti del \sigla{IIø}{7}{}}
\label{subsec:rivolti_ii7_minore}

\begin{regola}[Collocazione del \sigla{IIø}{7}{}]
	\label{reg:collocazione_ii7_minore}
	L'accordo semidiminuito si trova sul \textbf{II grado della scala minore armonica}; compare anche sul VI e sul VII grado della scala minore melodica ascendente.
\end{regola}

Oltre alla settima, anche la \textbf{quinta diminuita} è dissonante: entrambe devono scendere per grado congiunto in ogni collegamento.

\begin{regola}[Quinta diminuita a risoluzione obbligata]
	\label{reg:quinta_diminuita}
	Nella settima di terza specie sia la settima sia la \textbf{quinta diminuita} risolvono \textbf{per grado congiunto discendente} in tutti i collegamenti. La fondamentale è l'unica nota che sale (alla dominante).
\end{regola}

I rivolti del \sigla{IIø}{7}{} (\sigla{II}{6}{5}, \sigla{II}{4}{3}, \sigla{II}{2}{}) seguono le regole, i gradi al basso e i movimenti prescritti per la seconda specie (Capitolo \ref{chap:seconda_specie}), tenendo conto delle alterazioni del minore e dell'obbligo di far scendere la quinta diminuita.

\begin{center}
\begin{tabularx}{\linewidth}{lLLLL}
\toprule
\textbf{Stato} & \textbf{Cifratura} & \textbf{Grado al basso} & \textbf{Risolve su} & \textbf{Basso obbligato} \\
\midrule
Fond. & 7 & II & V (o \sigla{V}{7}{}) & no \\
I rv. & 6/5 & IV & V & sì (sale) \\
II rv. & 4/3 & VI & V & sì (scende) \\
III rv. & 4/2 & I (settima) & \sigla{V}{6}{} o \sigla{V}{6}{5} & sì (scende) \\
\bottomrule
\end{tabularx}
\end{center}

\begin{esempio}[\sigla{II}{7}{} $\rightarrow$ V in La minore: tutti gli stati]
	\label{ex:ii7_stati_minore}
	Stato fondamentale (completa e incompleta) e primo rivolto, poi secondo e terzo rivolto. La settima La scende a Sol$\sharp$, la quinta diminuita Fa scende a Mi; il Sol$\sharp$ (sensibile del V) non si raddoppia. Nel \sigla{II}{2}{} il basso La (la settima) scende al Sol$\sharp$ sul \sigla{V}{6}{}.
	\musica{es_cap13_semidim_la_1}
	\musica{es_cap13_semidim_la_2}
\end{esempio}

\begin{osservazione}[Omissione della quinta nel collegamento I $-$ \sigla{II}{7}{} $\rightarrow$ V]
	\label{oss:omissione_quinta}
	A differenza della seconda specie, nel collegamento I $-$ \sigla{II}{7}{} $\rightarrow$ V la quinta diminuita è un elemento caratteristico dell'accordo: alcuni teorici ne \textbf{vietano l'omissione}, altri accettano l'accordo senza quinta e con terza raddoppiata per evitare quinte parallele.
\end{osservazione}
